\documentclass[11pt]{article}
\usepackage[letterpaper,margin=1in]{geometry}
\usepackage[T1]{fontenc}
\usepackage{lmodern}
\usepackage{amsmath,amssymb,amsthm}
\usepackage[colorlinks=true,linkcolor=blue,urlcolor=blue,citecolor=blue]{hyperref}
\usepackage{microtype}

\newtheorem*{theorem}{Theorem}
\newcommand{\M}{\mathcal M}
\newcommand{\Ecal}{\mathcal E}

\title{Removing the Finite-Support Hypothesis:\\
Literature Resolution of the Question in arXiv:math/0004146}
\author{Run \texttt{fa\_banach\_001} --- literature-status note}
\date{August 11, 2026}

\begin{document}
\maketitle

\noindent\textbf{Status.} \texttt{literature\_already\_answered}. This note
identifies an exact later affirmative solution. It does not claim a new theorem.

\section*{The source question}
Let $(\M,\tau)$ be a semifinite von Neumann algebra and let $E$ be an
order-continuous symmetric Banach function space with the Fatou property.
Randrianantoanina's Proposition 2.6 in the arXiv draft (Proposition 3.4 in the
revised journal paper) considers a normalized sequence $(x_n)$ in
$E(\M,\tau)$ under the additional hypothesis
\begin{equation*}
 \text{there is }e\in\M\text{ with }e^2=e,\quad \tau(e)<\infty,
 \quad ex_n=x_n\quad(n\geq1).                         \tag{F}
\end{equation*}
It proves the following alternative: either there is $C>0$ such that
\begin{equation*}
 \int_0^1\left\|\sum_{i=1}^n r_i(t)a_i x_i\right\|_{E(\M,\tau)}dt
 \geq C\left(\sum_{i=1}^n|a_i|^2\right)^{1/2}        \tag{R}
\end{equation*}
for every $n$ and every scalar choice, or $(x_n)$ has a subsequence
equivalent to a disjointly supported sequence in $E$.

Immediately afterward the paper says: ``We do not know if condition (ii) can
be removed'' \cite[Remarks 3.5]{R}.  Condition (ii) is exactly (F).  The 2003
journal revision retains this sentence, so this was not merely an obsolete
line in the 2000 arXiv draft.

\section*{Exact later answer}
Huang and Sukochev prove the noncommutative Lindenstrauss--Tzafriri
dichotomy at the level of general operator bimodules \cite{HS}.  Their
published abstract states the result in the following form.

\begin{theorem}[Huang--Sukochev]
Let $p>0$ and let $\Ecal\subset L_p(\M,\tau)+\M$ be a noncommutative
quasi-Banach $\M$-bimodule with order-continuous quasi-norm.  For a bounded
(in particular, normalized) sequence $(x_i)$ in $\Ecal$, either there is a
constant $c>0$ for which
\[
 \int_0^1\left\|\sum_{i=1}^n r_i(t)a_i x_i\right\|_{\Ecal}dt
 \geq c\left(\sum_{i=1}^n|a_i|^2\right)^{1/2}
 \qquad(n=1,2,\ldots),
\]
for every scalar choice, or $(x_i)$ has a subsequence equivalent to a
sequence of disjoint elements.
\end{theorem}

Most decisively for provenance, the same abstract explicitly says that this
answers the question of Randrianantoanina in the 2003 paper \cite{HS}.

\section*{Why this removes precisely (F)}
Take $\Ecal=E(\M,\tau)$.  It is a Banach $\M$-bimodule, and order continuity
of $E$ gives order continuity of its norm.  As used already inside the source
proof, the symmetric space embeds continuously into
$L_1(\M,\tau)+\M$; hence the later theorem applies with $p=1$.  A normalized
sequence is bounded, and no common finite-trace support projection occurs in
the later hypotheses.

The first later alternative is word-for-word (R).  In the second alternative,
disjoint operator elements have orthogonal left and right supports.  The
standard singular-value disjointification used in the source's Proposition
2.4 identifies such a sequence, up to equivalence, with a disjointly
supported sequence in the underlying symmetric function space $E$.  Thus the
second alternative is also the source conclusion.

Consequently, condition (F) can be removed, affirmatively and without adding
the type-$2$ assumption used elsewhere in the 2003 paper.  In fact, the later
theorem is stronger: it works for order-continuous quasi-Banach operator
bimodules, not only symmetric Banach spaces with the Fatou property.

\section*{Scope and verification trail}
This packet addresses only Remarks 3.5/its arXiv-draft counterpart.  It does
not reclassify the paper's separate embedding theorems for noncommutative
Lorentz spaces.  The included 2003 publisher PDF verifies the final wording
of the source question.  The DOI metadata and the complete published abstract
of \cite{HS} are included as JSON snapshots; the abstract both states the
unrestricted alternative and names the source question as an application.

\begin{thebibliography}{9}
\bibitem{R}
N. Randrianantoanina,
\emph{Embeddings of $\ell_p$ into non-commutative spaces},
J. Aust. Math. Soc. \textbf{74} (2003), 331--350;
arXiv:math/0004146.

\bibitem{HS}
J. Huang and F. Sukochev,
\emph{Bounded sequences in noncommutative bimodules of
$\tau$-measurable operators},
Proc. Amer. Math. Soc. \textbf{154} (2026), no.~2, 793--806,
\href{https://doi.org/10.1090/proc/17465}{doi:10.1090/proc/17465}
(published online December 1, 2025).
\end{thebibliography}

\end{document}
